%% Elevark fra mattelab.no – laget automatisk av bygg.py. Rediger gjerne til eget bruk.
%% Kompiler med lualatex.
\documentclass[11pt,a4paper]{article}
\usepackage[norsk]{babel}
\usepackage[a4paper,margin=2.2cm,top=2cm,bottom=2.4cm]{geometry}
\usepackage{fontspec}
\usepackage{unicode-math}
\setmainfont{TeX Gyre Pagella}
\setsansfont{TeX Gyre Heros}[Scale=MatchLowercase]
\setmathfont{TeX Gyre Pagella Math}
\usepackage{xcolor}
\definecolor{blekk}{HTML}{16213A}
\definecolor{dempet}{HTML}{566077}
\definecolor{rute}{HTML}{CFC9B8}
\definecolor{m1}{HTML}{2F4AB0}
\definecolor{m2}{HTML}{99266D}
\usepackage[most]{tcolorbox}
\usepackage{qrcode}
\usepackage{needspace}
\usepackage{fancyhdr}
\usepackage[hidelinks]{hyperref}
\pdfvariable suppressoptionalinfo \numexpr 1+2+4+8+16+32+64+128+256+512\relax
\setlength{\parindent}{0pt}
\setlength{\parskip}{0.4em}
\pagestyle{fancy}
\fancyhf{}
\renewcommand{\headrulewidth}{0pt}
\fancyfoot[L]{\sffamily\footnotesize\color{dempet}mattelab.no · Brøk som areal · elevark uten fasit}
\fancyfoot[R]{\sffamily\footnotesize\color{dempet}Side \thepage}

\newcommand{\fase}[2]{\tcbox[on line, colframe=#2, colback=white, coltext=#2, boxrule=0.6pt, arc=2pt,
  boxsep=0pt, left=3pt, right=3pt, top=1.5pt, bottom=1.5pt, fontupper=\sffamily\bfseries\footnotesize]{#1}}
\newcommand{\steg}[4]{%
  \par\noindent{\color{rute}\rule{\linewidth}{0.4pt}}\par\nopagebreak
  \noindent\begin{minipage}[t]{2.1cm}\vspace{0pt}\fase{#1}{#2}\end{minipage}%
  \begin{minipage}[t]{\dimexpr\linewidth-2.1cm\relax}\vspace{0pt}#3\end{minipage}\par
  \vspace{#4}}
\newcommand{\oppgave}[2]{\par\Needspace{12\baselineskip}\vspace{0.6em}%
  {\fontsize{14}{17}\selectfont\bfseries{\color{m1}#1}\enspace #2}\par\vspace{0.2em}}

\begin{document}
\begin{minipage}[t]{0.72\linewidth}\vspace{0pt}
{\sffamily\small\color{dempet}Elevark · 1T · mattelab.no}\par\vspace{0.3em}
{\fontsize{24}{28}\selectfont\bfseries Brøk som areal}\par\vspace{1em}
{\sffamily\small Navn: \rule{4.6cm}{0.4pt}\quad Klasse: \rule{1.5cm}{0.4pt}\quad Dato: \rule{1.9cm}{0.4pt}}
\end{minipage}\hfill
\begin{minipage}[t]{0.22\linewidth}\vspace{0pt}\raggedleft
\qrcode[height=2.3cm]{https://mattelab.no/utforskere/brok-som-areal.html}\par\vspace{0.2em}
{\sffamily\scriptsize\color{dempet}Utforskeren}
\end{minipage}

\vspace{1em}
{\sffamily\small\color{dempet}Utforskeren ligger på \url{https://mattelab.no/utforskere/brok-som-areal.html}}\par
Hver oppgave har tre steg. \textbf{Gjett} først, uten å bruke utforskeren. \textbf{Test} så i utforskeren. \textbf{Forklar} til slutt med egne ord hvorfor det blir slik.

\oppgave{1}{Blir det mindre når vi ganger?}
\steg{Gjett}{m1}{Er $\frac{2}{3} \cdot \frac{3}{4}$ større eller mindre enn $\frac{2}{3}$? Gjett før du regner.}{2.2cm}
\steg{Test}{m2}{Vis $\frac{2}{3} \cdot \frac{3}{4}$. Tell de mørke rutene. Hvor mange små ruter er det i alt?}{2.6cm}
\steg{Forklar}{blekk}{Forklar med figuren hvorfor vi ganger tellerne med hverandre og nevnerne med hverandre.}{3.4cm}

\oppgave{2}{Hvorfor ikke $\frac{2}{5}$?}
\steg{Gjett}{m1}{En elev skriver $\frac{1}{2} + \frac{1}{3} = \frac{2}{5}$. Kan det stemme? Sammenlign $\frac{2}{5}$ med $\frac{1}{2}$.}{2.2cm}
\steg{Test}{m2}{Vis $\frac{1}{2} + \frac{1}{3}$. Hva skjer når begge brøkene deles i like store deler?}{2.6cm}
\steg{Forklar}{blekk}{Forklar hvorfor vi trenger felles nevner, og hva som er galt med å legge sammen tellere og nevnere hver for seg.}{3.4cm}

\oppgave{3}{Blir det større når vi deler?}
\steg{Gjett}{m1}{Hvor mange åttendeler er det i tre firedeler? Gjett hva $\frac{3}{4} : \frac{1}{8}$ er.}{2.2cm}
\steg{Test}{m2}{Vis $\frac{3}{4} : \frac{1}{8}$. Tell hvor mange ganger $\frac{1}{8}$ går opp.}{2.6cm}
\steg{Forklar}{blekk}{Forklar hvorfor svaret blir større enn $\frac{3}{4}$, og vis at regelen «gang med den omvendte brøken» gir samme svar.}{3.4cm}

\oppgave{4}{Når det ikke går opp}
\steg{Gjett}{m1}{Hvor mange ganger går $\frac{1}{4}$ opp i $\frac{2}{3}$? Er svaret et helt tall?}{2.2cm}
\steg{Test}{m2}{Vis $\frac{2}{3} : \frac{1}{4}$. Hvor mange hele ganger går det opp? Hvor stor er resten?}{2.6cm}
\steg{Forklar}{blekk}{Forklar hva svaret $2\frac{2}{3}$ betyr. Bruk ordene «hele ganger» og «rest».}{3.4cm}

\oppgave{5}{Når blir produktet større?}
\steg{Gjett}{m1}{Du ganger $\frac{3}{4}$ med et tall. Når blir svaret større enn $\frac{3}{4}$, og når blir det mindre?}{2.2cm}
\steg{Test}{m2}{Sammenlign $\frac{3}{4} \cdot \frac{1}{2}$ og $\frac{3}{4} \cdot \frac{3}{2}$. Prøv også å gange med $\frac{4}{4}$.}{2.6cm}
\steg{Forklar}{blekk}{Lag en regel for når et produkt blir større enn det første tallet. Forklar regelen med figuren.}{3.4cm}

\end{document}
